\section{A lower bound for point descendants}
\label{sec:descendants}

The spectral argument will require point descendants in a sequence of
large curve degrees.  Their virtual definition does not give positivity
directly.  We prove the required lower bound by a contraction that forgets
all components away from the marked component.  Its boundary image is too
small to contribute to the integral.

Throughout this section, \(X\) is a smooth projective complex Fano manifold
of dimension \(n>0\) and pseudoindex at least two.  Fix an embedding
\[
  \jmath:X\hookrightarrow\mathbb P^N,\qquad
  \jmath^*\mathcal O_{\mathbb P^N}(1)\cong\mathcal O_X(-aK_X),
\]
where \(a\) is a positive integer.  A morphism
\(f:\mathbb P^1\to X\) is \emph{free} if \(f^*T_X\) is globally generated.
For a numerical curve degree \(\beta\), set \(d_\beta=-K_X\cdot\beta\).
The notation \(\overline{\mathcal M}_{0,1}(X,\beta)\) includes the union
over actual curve classes with numerical degree \(\beta\).  This is a
finite union: maps of fixed projective degree belong to a finite-type
proper moduli stack, and numerical degree is locally constant.

We use rational Chow groups of Deligne--Mumford stacks.  In particular,
degrees of zero-dimensional stack cycles include their stabilizer
denominators.  The marked cotangent line has first Chern class \(\psi\).

\begin{proposition}[Point-descendant lower bound]
\label{prop:descendant-bound}
Suppose that \(\beta\) is represented by a nonconstant free map
\(\mathbb P^1\to X\).  Put \(d=d_\beta\) and \(b=ad\).  Then
\[
  \left\langle\tau_{d-2}(\mathrm{pt})\right\rangle_\beta
  \ge \frac{1}{b^{d-1}}.
\]
\end{proposition}

We will express this descendant as the degree of a nonempty effective
cycle in a weighted projective stack whose weights are at most \(b\).
The main task is to obtain that cycle from the virtual evaluation
fiber: irreducible maps will contribute with stack degree one, while
the boundary image will have too small a dimension to contribute.
After proving the bound, we show that free degrees are Zariski dense
in the numerical curve space.

\subsection{A free evaluation fiber}

There is a very general point \(x\in X\) such that every nonconstant
morphism \(\mathbb P^1\to X\) whose image contains \(x\) is free.
Here is the argument, including the quantifier over special maps.
Stratify the nonfree locus in each morphism scheme into smooth
irreducible locally closed reduced strata.  There are countably many
such strata.  If the evaluation from \(\mathbb P^1\times T\), for one
stratum \(T\), were dominant, generic smoothness would give a point
\((p,f)\) where its differential is surjective.  Its image is contained
in the image of
\[
  H^0(\mathbb P^1,f^*T_X)\longrightarrow(f^*T_X)_p.
\]
Indeed, moving \(p\) gives directions also obtained from
reparametrizations of the source.  Splitting
\(f^*T_X=\bigoplus_j\mathcal O(a_j)\), surjectivity at one point implies
\(a_j\ge0\) for every \(j\), contrary to the choice of \(T\).
Thus each such evaluation image has proper closure, and their
complement gives the assertion.

Evaluation is smooth at a free map: the obstruction groups
\(H^1(\mathbb P^1,f^*T_X)\) and
\(H^1(\mathbb P^1,f^*T_X(-p))\) vanish.  Consequently the degree
\(\beta\) in Proposition~\ref{prop:descendant-bound} has maps
through an open subset of \(X\).  Choose \(x\) in this open subset and
outside all the exceptional closed subsets just considered.

Let
\[
  F=\operatorname{ev}^{-1}(x)
    \subset\overline{\mathcal M}_{0,1}(X,\beta),\qquad h=d-2.
\]
The fiber carries the refined virtual class
\[
  [F]^{\mathrm{vir}}
  =x^![\overline{\mathcal M}_{0,1}(X,\beta)]^{\mathrm{vir}}
  \in A_h(F)_{\mathbb Q}.
\]
Let \(U\subset F\) be the open substack with irreducible source.
It is nonempty, smooth of dimension \(h\), and
\[
  [F]^{\mathrm{vir}}|_U=[U].
\]
Only this open part is asserted to be unobstructed.

To use this equality of classes on \(U\), we will construct a proper
morphism that preserves the marked cotangent class, has generic stack
degree one on \(U\), and sends \(F\setminus U\) to a locus of
dimension less than \(h\).  These properties will make the
pushforward of \([F]^{\mathrm{vir}}\) effective.

\subsection{Contracting away from the marking}

Compose maps with \(\jmath\).  Their projective degree is \(b\), and
this composition changes no stable domain, since it contracts no
nonconstant component.  At a geometric point, retain the component
carrying the marking and replace each attached tree by a common zero
of the map sections at its attachment point.  More precisely, multiply
the original sections restricted to the marked component by a common
factor whose vanishing order at each attachment is the projective
degree of the removed tree.  The resulting tuple has total degree
\(b\) and no common zero at the marking.  This description also
allows the marked component to be constant.

We now construct these data in families on
\(\overline{\mathcal M}_{0,1}(\mathbb P^N,b)\) and track the marked
cotangent line.  The contraction also follows from the one-marked
intermediate-space construction of
Mustaţă--Mustaţă~\cite[Definition~1.2 and Proposition~1.3]
{MustataMustata2007}, with the fine-stack construction of
\cite[Proposition~1.7]{MustataMustataChow}.

Write \(\pi:\mathcal C\to\overline{\mathcal M}_{0,1}(\mathbb P^N,b)\)
for the universal curve, \(s\) for its marking, and
\(\mathcal L\) for the pullback of \(\mathcal O_{\mathbb P^N}(1)\).
The moduli stack is smooth and its boundary has the usual independent
node-smoothing parameters: on a genus-zero tree,
\(H^1(C,f^*T_{\mathbb P^N})=0\).  Thus the universal curve is regular,
with local nodal models \(uv=t\).

For each boundary branch separating an unmarked tail of total degree
\(e\), take in a local node-smoothing chart the Cartier divisor
consisting of its unmarked side.  The marked side and degree \(e\)
are intrinsic, so summing these local divisors over all such branches
is invariant under permutations of equal-degree branches and glues
in families.  Let \(\Delta_e\) be this sum, counting repeated branches
once each.
Only \(1\le e\le b-1\) occurs.  A tail of degree \(b\) would leave a
degree-zero marked side with only one external attachment; a finite
tree of contracted components with one marking and one external
attachment cannot satisfy stability.

Set
\[
  \mathcal L^+
    =\mathcal L\Bigl(\sum_{e=1}^{b-1}e\Delta_e\Bigr).
\]
On a nodal fiber, twisting by an unmarked-side divisor transfers one
unit of degree across that node toward the marked side.  To compute
the resulting multidegree, root the dual tree at the marked component.
For a nonroot component \(v\), let \(b_v\) be its original degree
and \(E_v\) the total degree of the subtree rooted at \(v\).
If \(w\) runs over its children, then
\[
  \deg_v\mathcal L^+=b_v-E_v+\sum_w E_w=0.
\]
On the root, of original degree \(b_0\), the degree is
\(b_0+\sum_wE_w=b\).  Thus \(\mathcal L^+\) has degree \(b\)
on the component containing \(s\), and degree zero elsewhere.

The line bundle \(\mathcal O_{\mathcal C}(s)\) has degree one on that
marked component and degree zero elsewhere.  On each fiber it is
globally generated with \(h^0=2\) and \(h^1=0\).  Cohomology and base
change consequently give a rank-two bundle
\(\pi_*\mathcal O_{\mathcal C}(s)\), and its evaluation map contracts
the universal curve onto a \(\mathbb P^1\)-bundle.  This contraction
is an isomorphism on the marked component and near the section.

The line bundle \(\mathcal L^+(-bs)\) has multidegree zero on every
fiber.  Its fiberwise \(h^0=1\), \(h^1=0\), and evaluation map show
that it is pulled back from the base.  Hence \(\mathcal L^+\)
descends to a relative degree-\(b\) line bundle on the contracted
\(\mathbb P^1\)-bundle.  Fiberwise pullback identifies its sections
with those of \(\mathcal L^+\): a section is constant on every
contracted tail and is determined by its value on the marked component.
The corresponding \(H^0\) bundles agree by base change, so the
inclusion \(\mathcal L\to\mathcal L^+\) sends the original map
sections to sections of the descended bundle in families.  Restricted
to the marked component, this inclusion multiplies them by a section
vanishing to order \(e\) at the attachment of a tree of degree \(e\).
The original map sections have no common zero, so the new common zero
has exactly that multiplicity.  These constructions commute with base
change and produce the retained data in families.  Since the
contraction is an isomorphism near the marking, it preserves the
marked cotangent line.

\subsection{The weighted projective fiber}

Choose target coordinates with \(x=[1:0:\cdots:0]\).  Locally on
the base, identify the contracted bundle with a product
\(\mathbb P^1\)-bundle and the descended line bundle with
\(\mathcal O_{\mathbb P^1}(b)\).  Put the mark at infinity and
normalize the degree-\(b\) leading coefficient vector of the
polynomial sections to \((1,0,\ldots,0)\).  Translation of the
affine coordinate uniquely eliminates the coefficient of degree
\(b-1\) in the first polynomial.  After these normalizations, two
local choices differ only by scaling the affine coordinate.
Consequently the coefficient tuples glue as points of the quotient
stack for this residual scaling action.

The remaining coefficients have weights
\[
  \underbrace{1,\ldots,1}_{N},\quad
  \underbrace{2,\ldots,2}_{N+1},\quad\ldots,\quad
  \underbrace{b,\ldots,b}_{N+1}.
\]
They are not all zero.  Otherwise the polynomial tuple would be
\((z^b,0,\ldots,0)\), a constant map with all degree concentrated at
one basepoint.  A constant marked component has at least two distinct
attachments by stability, so cannot produce that tuple; a
nonconstant marked component cannot produce it either.

Let \(w_1,\ldots,w_M\) be the displayed list of weights and put
\[
 \mathcal W=[(\mathbb A^M\setminus\{0\})/\mathbb G_m],
 \qquad \lambda\cdot(x_1,\ldots,x_M)
       =(\lambda^{w_1}x_1,\ldots,\lambda^{w_M}x_M).
\]
This is the weighted projective stack in question. We have constructed
a proper morphism
\[
  \Phi:F\longrightarrow\mathcal W.
\]
The map is proper because \(F\) is proper and \(\mathcal W\) is
separated. For the reparametrization \(z\mapsto\lambda z\), the
normalized coefficient of \(z^{b-j}\) has weight \(-j\), and the
cotangent \(d(1/z)\) has weight \(-1\). Replacing \(\lambda\) by
\(\lambda^{-1}\) gives the displayed positive weights and identifies
the cotangent line with the pullback of
\(\mathcal O_{\mathcal W}(1)\), as in the one-marked weighted construction
of~\cite[proof of Lemma~3.3]{MustataMustata2007}.  Thus, with
\(\xi=c_1(\mathcal O_{\mathcal W}(1))\),
\[
  \psi=\Phi^*\xi.
\]

On the open substack \(U\), the tuple has no basepoints and recovers
the entire pointed map.  Its image is the locus of basepoint-free
tuples whose maps factor through \(\jmath(X)\) and have numerical
degree \(\beta\) there.  A residual scaling fixes the tuple exactly
when the associated reparametrization preserves the pointed map, so
the stabilizers are the automorphism groups of the pointed maps.  We
therefore have an identification of stacks onto this locus, with
generic stack degree one on every irreducible component of \(U\).

\subsection{The boundary does not contribute}

Consider a boundary stratum of \(F\).  Let \(d_0\) be the
anticanonical degree of the marked component and let \(u\ge1\) be
the number of trees directly attached to it.  Its image in
\(\mathcal W\) is determined by the marked map, the mark and
attachment points, and the degrees of those trees.  It does not
retain the maps on the trees: a tree of anticanonical degree \(d_i\)
is recorded only by its common-zero multiplicity \(ad_i\).

If \(d_0>0\), the marked map is free, since it passes through \(x\).
The space of its parametrized maps has dimension \(n+d_0\).
Adding \(u+1\) marked points, imposing evaluation \(x\), and dividing
by source reparametrization gives dimension
\[
  n+d_0+(u+1)-n-3=d_0+u-2.
\]
If \(d_0=0\), the map is the constant \(x\), stability gives \(u\ge2\),
and the remaining data have dimension
\(\dim\mathcal M_{0,u+1}=u-2\), the same formula.

Every attached tree contains a nonconstant rational component.
The pseudoindex assumption implies
\[
  d-d_0\ge2u,\qquad
  d_0+u-2\le d-u-2\le h-1.
\]
There are finitely many discrete boundary types.  We conclude that
\[
  \dim\Phi(F\setminus U)<h.
\]
This estimate concerns the image only; it makes no smoothness or
dimension assertion about the spaces of tail maps.

Write \(Z_\alpha\) for the closures in \(F\) of the irreducible
components of \(U\).  Chow localization and
\([F]^{\mathrm{vir}}|_U=[U]\) express
\[
  [F]^{\mathrm{vir}}-\sum_\alpha[Z_\alpha]
\]
as a class supported on \(F\setminus U\).  Its pushforward to
\(\mathcal W\) vanishes, since every \(h\)-dimensional boundary cycle
has image of smaller dimension.  The generic degree-one statement
on \(U\) therefore gives
\[
  \Phi_*[F]^{\mathrm{vir}}
    =\sum_\alpha[\overline{\Phi(U_\alpha)}].
\]
This is a nonempty effective sum of \(h\)-dimensional integral
stack cycles.  Virtual effectiveness has been established only
after pushforward; no effectiveness of \([F]^{\mathrm{vir}}\)
itself was assumed.

\subsection{Degrees in a weighted projective stack}

\begin{lemma}
\label{lem:weighted-degree-lower}
Let \(\mathcal P\) be a weighted projective stack with positive
integer weights at most \(b\), and let \(Z\subset\mathcal P\) be an
integral closed substack of dimension \(h\).  Then
\[
  \int_Z c_1(\mathcal O_{\mathcal P}(1))^h
  \ge b^{-(h+1)}.
\]
\end{lemma}

\begin{proof}
The affine cone of \(Z\) is defined by an ideal homogeneous for the
weighted grading.  Degenerate it to a monomial initial ideal.
This is a flat grading-preserving degeneration.  Removing the
origin and quotienting by the grading action of \(\mathbb G_m\)
gives a flat family of closed substacks of the fixed proper
weighted projective stack.  Its \(\mathcal O(1)\)-degree is constant.

The top-dimensional cycle of the monomial limit is a nonempty sum
of coordinate weighted projective substacks, with positive integer
multiplicities.  A component of dimension \(h\) has the form
\(\mathbb P(w_0,\ldots,w_h)\), with \(1\le w_j\le b\), and
\[
  \int_{\mathbb P(w_0,\ldots,w_h)}
       c_1(\mathcal O(1))^h
   =\frac1{w_0\cdots w_h}.
\]
For clarity, the power-coordinate map
\[
  \mathbb P^h\longrightarrow\mathbb P(w_0,\ldots,w_h),
  \qquad [x_0:\cdots:x_h]\longmapsto
          [x_0^{w_0}:\cdots:x_h^{w_h}]
\]
is the map of stacks specified by the line bundle
\(\mathcal O_{\mathbb P^h}(1)\) and these sections of its powers.
It is representable and finite, pulls back \(\mathcal O(1)\) to
\(\mathcal O_{\mathbb P^h}(1)\), and has stack degree
\(\prod_jw_j\).  On the dense tori the coarse degree is
\(\prod_jw_j/g\), where \(g=\gcd(w_0,\ldots,w_h)\); the target's
generic stabilizer has order \(g\), giving the stated stack degree.
This also covers \(h=0\), when the target is \(B\mu_{w_0}\).
The displayed integral follows by the projection formula.

Each component degree is at least \(b^{-(h+1)}\), so the same lower
bound holds for their nonempty positive sum and hence for \(Z\).
\end{proof}

\begin{proof}[Proof of Proposition~\ref{prop:descendant-bound}]
By the definition of the point descendant and by the projection formula,
\[
  \left\langle\tau_{d-2}(\mathrm{pt})\right\rangle_\beta
  =\int_{[F]^{\mathrm{vir}}}\psi^h
  =\int_{\Phi_*[F]^{\mathrm{vir}}}\xi^h.
\]
The last cycle is the nonempty effective sum obtained above.
All weights of \(\mathcal W\) are at most \(b\).
Lemma~\ref{lem:weighted-degree-lower} now gives the bound
\(b^{-(h+1)}=b^{-(d-1)}\).
\end{proof}

\subsection{Enough free degrees}

The next lemma lets us apply the descendant bound in a direction
avoiding any specified nonzero polynomial relation.  Write
\(N_1(X)_{\mathbb C}\) for the complex vector space of numerical
curve classes.

\begin{lemma}
\label{lem:free-dense}
The numerical degrees represented by nonconstant free rational maps
are Zariski dense in \(N_1(X)_{\mathbb C}\).  They are closed under
addition and under multiplication by positive integers.
\end{lemma}

\begin{proof}
We first prove closure under addition.  Take two free maps of degrees
\(\beta_1,\beta_2\).  Their evaluation images contain nonempty open
subsets of \(X\).  Deform them so that they meet at a common point
and glue there.  On the resulting two-component genus-zero curve
\(C\), let \(f:C\to X\) be the glued map and choose a smooth point
\(p\) of the first component.  Normalization gives
\[
  H^1(C,f^*T_X)=H^1(C,f^*T_X(-p))=0.
\]
Indeed the corresponding cohomology groups vanish on both
components, and sections on the second component evaluate
surjectively at the node.  Deformation theory therefore permits
a pointed smoothing of the node.  Semicontinuity of the second
vanishing shows that the smoothed map is free.  Its degree is
\(\beta_1+\beta_2\).

We next show that free degrees span the numerical vector space.
Smooth complex Fano manifolds are rationally chain connected by
\cite[Corollary~3.2]{Campana1992}, hence rationally connected by
\cite[Theorem~2.49]{Debarre2011}; in characteristic zero this is
equivalent to separable rational connectedness
\cite[Definition~2.20 and the following discussion]{Debarre2011}.
The rational-connectedness results are due to Campana and
Koll\'ar--Miyaoka--Mori~\cite{Campana1992,KollarMiyaokaMori1992}.
By \cite[Theorem~1.3]{TianZong2014}, the Chow group of one-cycles on \(X\)
is generated by rational curves.  In particular, rational curve
classes span \(N_1(X)_{\mathbb C}\).
A map is \emph{very free}
if its pulled-back tangent bundle is ample.  For a morphism
\(f_0:\mathbb P^1\to X\), comb smoothing attaches a nonempty finite
collection of very free rational maps as teeth and deforms the resulting
tree to a free rational map.  More precisely,
Tian--Zong~\cite[Proposition~2.4]{TianZong2014}, with handle
\(\mathbb P^1\) and line-bundle twist of degree \(-1\), gives a
smoothing \(f\) with \(H^1(\mathbb P^1,f^*T_X(-1))=0\).
Splitting on \(\mathbb P^1\) shows that \(f\) is free.  Its numerical
degree is the sum of the handle degree and the tooth degrees.
The tooth degrees are free, and their sum is free by closure under
addition.  Every rational degree is consequently a difference of
free degrees, which proves the spanning assertion.

Choose \(r=\dim N_1(X)_{\mathbb C}\) linearly independent free
degrees from the spanning set.  Every positive integral combination
of these degrees is free by the preceding addition argument.
Such combinations are Zariski dense: in the basis formed by the
chosen degrees, a polynomial vanishing on
\(\mathbb Z_{>0}^r\) vanishes by successive application of the
one-variable root bound.  Finally, covers of a free map remain
free, since pullback preserves the nonnegative summands of
its split tangent bundle.  This gives the assertion about positive
integer multiples directly as well.
\end{proof}
